% =====================================================================
% FONTES DEPENDENTES --- ANALISE DE MALHAS
% 10 questoes visuais. Com 7 questoes dependentes no banco nodal visual,
% completa o conjunto de 17 solicitado (a questao original + 16 novas).
% =====================================================================

\subsubsection*{Banco especial --- Fontes dependentes em análise de malhas}

\begin{atencao}[Leitura obrigatória do circuito]
As questões abaixo usam topologias distintas: fontes controladas por corrente e
por tensão, ramos externos e compartilhados, supermalha, três malhas,
determinação de ganho e potência. Adote as correntes de malha indicadas no
sentido horário.
\end{atencao}

\blocoaprofundamento

\begin{questao}[3]{}
A fonte de tensão dependente vale $2{,}5I_1$ volts e impulsiona a malha 2. Determine $I_1$, $I_2$ e a corrente real no resistor central.
\circuitoq{\begin{circuitikz}[american,scale=.88,transform shape]
\draw (0,0) to[battery1,l_={$\SI{20}{\volt}$}] (0,3) to[R,l={$\SI{5}{\ohm}$}] (3,3) coordinate(t);
\draw (t) to[R,l={$\SI{5}{\ohm}$}] (3,0) coordinate(b)--(0,0);
\draw (t) to[R,l={$\SI{5}{\ohm}$}] (6,3) to[cvsource,l={$2{,}5I_1$}] (6,0)--(b);
\draw[->,loopcol,line width=1.2pt] (1.4,1.5) arc (0:310:.48); \node[loopcol] at(1.4,1.5){\footnotesize$I_1$};
\draw[->,loopcol,line width=1.2pt] (4.6,1.5) arc (0:310:.48); \node[loopcol] at(4.6,1.5){\footnotesize$I_2$};
\end{circuitikz}}
\resposta{$I_1=\SI{3.2}{\ampere}$; $I_2=\SI{2.4}{\ampere}$; $I_c=\SI{0.8}{\ampere}$}
\resolucao{Malha 1: $10I_1-5I_2=20$. Malha 2: $-5I_1+10I_2=2{,}5I_1$, ou $-7{,}5I_1+10I_2=0$. Daí $I_1=3{,}2$ A e $I_2=2{,}4$ A. No resistor comum, $I_c=I_1-I_2=0{,}8$ A.}
\end{questao}

\begin{questao}[3]{}
A fonte dependente vale $4i_x$ volts, com $i_x=I_1-I_2$ no resistor central. Determine as correntes de malha.
\circuitoq{\begin{circuitikz}[american,scale=.88,transform shape]
\draw (0,0) to[battery1,l_={$\SI{24}{\volt}$}] (0,3) to[R,l={$\SI{4}{\ohm}$}] (3,3) coordinate(t);
\draw (t) to[R,l={$\SI{4}{\ohm}$},i>^={$i_x$}] (3,0) coordinate(b)--(0,0);
\draw (t) to[R,l={$\SI{8}{\ohm}$}] (6,3) to[cvsource,l={$4i_x$}] (6,0)--(b);
\draw[->,loopcol,line width=1.2pt] (1.4,1.5) arc (0:310:.48); \node[loopcol] at(1.4,1.5){\footnotesize$I_1$};
\draw[->,loopcol,line width=1.2pt] (4.6,1.5) arc (0:310:.48); \node[loopcol] at(4.6,1.5){\footnotesize$I_2$};
\end{circuitikz}}
\resposta{$I_1=\SI{4}{\ampere}$; $I_2=\SI{2}{\ampere}$}
\resolucao{Malha 1: $8I_1-4I_2=24$. Na malha 2, $-4I_1+12I_2=4(I_1-I_2)$, portanto $-8I_1+16I_2=0$. Assim $I_1=4$ A e $I_2=2$ A.}
\end{questao}

\begin{questao}[3]{}
A tensão de controle $v_x$ é a tensão no resistor de $\SI{4}{\ohm}$ da malha 1. A fonte dependente vale $0{,}5v_x$. Determine $I_1$ e $I_2$.
\circuitoq{\begin{circuitikz}[american,scale=.88,transform shape]
\draw (0,0) to[battery1,l_={$\SI{24}{\volt}$}] (0,3) to[R,l={$\SI{4}{\ohm}$},v^>={$v_x$}] (3,3) coordinate(t);
\draw (t) to[R,l={$\SI{4}{\ohm}$}] (3,0) coordinate(b)--(0,0);
\draw (t) to[R,l={$\SI{4}{\ohm}$}] (6,3) to[cvsource,l={$0{,}5v_x$}] (6,0)--(b);
\draw[->,loopcol,line width=1.2pt] (1.4,1.5) arc (0:310:.48); \node[loopcol] at(1.4,1.5){\footnotesize$I_1$};
\draw[->,loopcol,line width=1.2pt] (4.6,1.5) arc (0:310:.48); \node[loopcol] at(4.6,1.5){\footnotesize$I_2$};
\end{circuitikz}}
\resposta{$I_1=\SI{4.8}{\ampere}$; $I_2=\SI{3.6}{\ampere}$}
\resolucao{$v_x=4I_1$, logo a fonte vale $2I_1$ V. As equações são $8I_1-4I_2=24$ e $-4I_1+8I_2=2I_1$. Da segunda, $I_2=0{,}75I_1$; substituindo na primeira, $I_1=4{,}8$ A e $I_2=3{,}6$ A.}
\end{questao}

\begin{questao}[3]{}
A fonte de corrente dependente está na periferia da malha 2 e impõe $I_2=0{,}5I_1$. Determine $I_1$, $I_2$ e o módulo da tensão na fonte de corrente.
\circuitoq{\begin{circuitikz}[american,scale=.88,transform shape]
\draw (0,0) to[battery1,l_={$\SI{30}{\volt}$}] (0,3) to[R,l={$\SI{5}{\ohm}$}] (3,3) coordinate(t);
\draw (t) to[R,l={$\SI{5}{\ohm}$}] (3,0) coordinate(b)--(0,0);
\draw (t) to[R,l={$\SI{10}{\ohm}$}] (6,3) to[cisource,l={$0{,}5I_1$},invert] (6,0)--(b);
\draw[->,loopcol,line width=1.2pt] (1.4,1.5) arc (0:310:.48); \node[loopcol] at(1.4,1.5){\footnotesize$I_1$};
\draw[->,loopcol,line width=1.2pt] (4.6,1.5) arc (0:310:.48); \node[loopcol] at(4.6,1.5){\footnotesize$I_2$};
\end{circuitikz}}
\resposta{$I_1=\SI{4}{\ampere}$; $I_2=\SI{2}{\ampere}$; $|v_s|=\SI{10}{\volt}$}
\resolucao{A fonte na periferia fixa diretamente $I_2=0{,}5I_1$. Na malha 1: $10I_1-5I_2=30$. Portanto $7{,}5I_1=30$, resultando $I_1=4$ A e $I_2=2$ A. Na malha 2, as quedas resistivas somam $10(2)+5(2-4)=10$ V em módulo; a fonte sustenta essa tensão.}
\end{questao}

\begin{questao}[3]{}
A fonte de corrente dependente está no ramo comum. Sua corrente, de baixo para cima, é $0{,}1v_x$, onde $v_x$ é a tensão no resistor esquerdo de $\SI{5}{\ohm}$. Resolva pela supermalha.
\circuitoq{\begin{circuitikz}[american,scale=.88,transform shape]
\draw (0,0) to[battery1,l_={$\SI{40}{\volt}$}] (0,3) to[R,l={$\SI{5}{\ohm}$},v^>={$v_x$}] (3,3) coordinate(t);
\draw (t) to[cisource,l={$0{,}1v_x$},invert] (3,0) coordinate(b)--(0,0);
\draw (t) to[R,l={$\SI{5}{\ohm}$}] (6,3)--(6,0)--(b);
\draw[->,loopcol,line width=1.2pt] (1.4,1.5) arc (0:310:.48); \node[loopcol] at(1.4,1.5){\footnotesize$I_1$};
\draw[->,loopcol,line width=1.2pt] (4.6,1.5) arc (0:310:.48); \node[loopcol] at(4.6,1.5){\footnotesize$I_2$};
\end{circuitikz}}
\resposta{$I_1=\SI{3.2}{\ampere}$; $I_2=\SI{4.8}{\ampere}$}
\resolucao{Como $v_x=5I_1$, a restrição é $I_2-I_1=0{,}1(5I_1)=0{,}5I_1$, logo $I_2=1{,}5I_1$. Na supermalha: $40=5I_1+5I_2$. Daí $I_1=3{,}2$ A e $I_2=4{,}8$ A.}
\end{questao}

\begin{questao}[3]{}
A fonte de corrente dependente do ramo comum fornece, de baixo para cima, $2i_x$, sendo $i_x=I_1$ no resistor esquerdo. Determine as correntes por supermalha.
\circuitoq{\begin{circuitikz}[american,scale=.88,transform shape]
\draw (0,0) to[battery1,l_={$\SI{40}{\volt}$}] (0,3) to[R,l={$\SI{5}{\ohm}$},i>^={$i_x$}] (3,3) coordinate(t);
\draw (t) to[cisource,l={$2i_x$},invert] (3,0) coordinate(b)--(0,0);
\draw (t) to[R,l={$\SI{5}{\ohm}$}] (6,3)--(6,0)--(b);
\draw[->,loopcol,line width=1.2pt] (1.4,1.5) arc (0:310:.48); \node[loopcol] at(1.4,1.5){\footnotesize$I_1$};
\draw[->,loopcol,line width=1.2pt] (4.6,1.5) arc (0:310:.48); \node[loopcol] at(4.6,1.5){\footnotesize$I_2$};
\end{circuitikz}}
\resposta{$I_1=\SI{2}{\ampere}$; $I_2=\SI{6}{\ampere}$}
\resolucao{A restrição é $I_2-I_1=2I_1$, portanto $I_2=3I_1$. A LKT da supermalha é $40=5I_1+5I_2$. Substituindo, $40=20I_1$, daí $I_1=2$ A e $I_2=6$ A.}
\end{questao}

\blocodesafio

\begin{questao}[4]{}
No circuito de três malhas, a fonte dependente da terceira malha vale $2{,}5I_1$ volts e impulsiona $I_3$. Determine $I_1$, $I_2$ e $I_3$.
\circuitoq{\begin{circuitikz}[american,scale=.70,transform shape]
\draw (0,0) to[battery1,l_={$\SI{20}{\volt}$}] (0,3) to[R,l={$\SI{5}{\ohm}$}] (2.6,3) coordinate(t1);
\draw (t1) to[R,l={$\SI{5}{\ohm}$}] (2.6,0) coordinate(b1)--(0,0);
\draw (t1) to[R,l={$\SI{5}{\ohm}$}] (5.2,3) coordinate(t2);
\draw (t2) to[R,l={$\SI{5}{\ohm}$}] (5.2,0) coordinate(b2)--(b1);
\draw (t2) to[R,l={$\SI{5}{\ohm}$}] (7.8,3) to[cvsource,l={$2{,}5I_1$}] (7.8,0)--(b2);
\draw[->,loopcol] (1.2,1.5) arc (0:310:.4); \node[loopcol] at(1.2,1.5){\scriptsize$I_1$};
\draw[->,loopcol] (3.9,1.5) arc (0:310:.4); \node[loopcol] at(3.9,1.5){\scriptsize$I_2$};
\draw[->,loopcol] (6.6,1.5) arc (0:310:.4); \node[loopcol] at(6.6,1.5){\scriptsize$I_3$};
\end{circuitikz}}
\resposta{$I_1=\dfrac{8}{3}\,\si{\ampere}$; $I_2=I_3=\dfrac{4}{3}\,\si{\ampere}$}
\resolucao{O sistema é $10I_1-5I_2=20$, $-5I_1+15I_2-5I_3=0$ e $-5I_2+10I_3=2{,}5I_1$. Resolvendo, $I_1=8/3$ A e $I_2=I_3=4/3$ A.}
\end{questao}

\begin{questao}[4]{}
A fonte de corrente dependente está entre as malhas 2 e 3 e impõe $I_3-I_2=0{,}5I_1$. Determine as três correntes usando uma supermalha entre 2 e 3.
\circuitoq{\begin{circuitikz}[american,scale=.70,transform shape]
\draw (0,0) to[battery1,l_={$\SI{30}{\volt}$}] (0,3) to[R,l={$\SI{5}{\ohm}$}] (2.6,3) coordinate(t1);
\draw (t1) to[R,l={$\SI{5}{\ohm}$}] (2.6,0) coordinate(b1)--(0,0);
\draw (t1) to[R,l={$\SI{5}{\ohm}$}] (5.2,3) coordinate(t2);
\draw (t2) to[cisource,l={$0{,}5I_1$},invert] (5.2,0) coordinate(b2)--(b1);
\draw (t2) to[R,l={$\SI{10}{\ohm}$}] (7.8,3)--(7.8,0)--(b2);
\draw[->,loopcol] (1.2,1.5) arc (0:310:.4); \node[loopcol] at(1.2,1.5){\scriptsize$I_1$};
\draw[->,loopcol] (3.9,1.5) arc (0:310:.4); \node[loopcol] at(3.9,1.5){\scriptsize$I_2$};
\draw[->,loopcol] (6.6,1.5) arc (0:310:.4); \node[loopcol] at(6.6,1.5){\scriptsize$I_3$};
\end{circuitikz}}
\resposta{$I_1=\SI{3}{\ampere}$; $I_2=0$; $I_3=\SI{1.5}{\ampere}$}
\resolucao{Malha 1: $10I_1-5I_2=30$. Supermalha 2--3: $-5I_1+10I_2+10I_3=0$. Restrição: $I_3-I_2=0{,}5I_1$. O sistema fornece $I_1=3$ A, $I_2=0$ e $I_3=1{,}5$ A.}
\end{questao}

\begin{questao}[4]{}
A fonte dependente vale $kI_1$ volts. Determine $k$ para que $I_2=\SI{2}{\ampere}$ e calcule $I_1$.
\circuitoq{\begin{circuitikz}[american,scale=.88,transform shape]
\draw (0,0) to[battery1,l_={$\SI{20}{\volt}$}] (0,3) to[R,l={$\SI{5}{\ohm}$}] (3,3) coordinate(t);
\draw (t) to[R,l={$\SI{5}{\ohm}$}] (3,0) coordinate(b)--(0,0);
\draw (t) to[R,l={$\SI{5}{\ohm}$}] (6,3) to[cvsource,l={$kI_1$}] (6,0)--(b);
\draw[->,loopcol,line width=1.2pt] (1.4,1.5) arc (0:310:.48); \node[loopcol] at(1.4,1.5){\footnotesize$I_1$};
\draw[->,loopcol,line width=1.2pt] (4.6,1.5) arc (0:310:.48); \node[loopcol] at(4.6,1.5){\footnotesize$I_2$};
\end{circuitikz}}
\resposta{$I_1=\SI{3}{\ampere}$; $k=\dfrac{5}{3}\,\ohm\approx\SI{1.67}{\ohm}$}
\resolucao{A malha 1 dá $10I_1-5I_2=20$. Impondo $I_2=2$ A, resulta $I_1=3$ A. Na malha 2, $-5I_1+10I_2=kI_1$. Logo $-15+20=3k$, portanto $k=5/3\,\ohm$.}
\end{questao}

\begin{questao}[4]{}
A fonte de tensão dependente está no ramo comum e possui tensão, de cima para baixo, $v_d=2I_1$. Determine $I_1$, $I_2$, a corrente real na fonte dependente e sua potência.
\circuitoq{\begin{circuitikz}[american,scale=.88,transform shape]
\draw (0,0) to[battery1,l_={$\SI{20}{\volt}$}] (0,3) to[R,l={$\SI{4}{\ohm}$}] (3,3) coordinate(t);
\draw (t) to[cvsource,l={$v_d=2I_1$}] (3,0) coordinate(b)--(0,0);
\draw (t) to[R,l={$\SI{4}{\ohm}$}] (6,3)--(6,0)--(b);
\draw[->,loopcol,line width=1.2pt] (1.4,1.5) arc (0:310:.48); \node[loopcol] at(1.4,1.5){\footnotesize$I_1$};
\draw[->,loopcol,line width=1.2pt] (4.6,1.5) arc (0:310:.48); \node[loopcol] at(4.6,1.5){\footnotesize$I_2$};
\end{circuitikz}}
\resposta{$I_1=\dfrac{10}{3}\,\si{\ampere}$; $I_2=\dfrac{5}{3}\,\si{\ampere}$; $i_d=\dfrac{5}{3}\,\si{\ampere}$ para baixo; $P_d=\dfrac{100}{9}\,\si{\watt}\approx\SI{11.11}{\watt}$ absorvidos}
\resolucao{Na malha 1, $4I_1+v_d=20$ e $v_d=2I_1$, então $I_1=10/3$ A. Na malha 2, $4I_2-v_d=0$, logo $I_2=5/3$ A. A corrente real no ramo comum é $i_d=I_1-I_2=5/3$ A, para baixo. Como $v_d=20/3$ V e a corrente entra no terminal positivo superior, a fonte absorve $P=v_di_d=100/9$ W.}
\end{questao}

\gabarito[Fontes dependentes --- análise de malhas]
